% ========================================================
%  习近平新时代中国特色社会主义思想概论  期末模拟试卷（二）（含参考答案）
%  覆盖范围：导论、第一章至第十二章（除第九章）
% ========================================================
\documentclass[12pt, a4paper]{article}

% ---- 中文支持 ----
\usepackage[UTF8]{ctex}

% ---- 数学 ----
\usepackage{amsmath, amssymb, amsthm}
\usepackage{mathtools}

% ---- 页面布局 ----
\usepackage[top=2.5cm, bottom=2.5cm, left=2.8cm, right=2.8cm]{geometry}

% ---- 表格 ----
\usepackage{booktabs, array, multirow, makecell}
\usepackage{tabularx}

% ---- 页眉页脚 ----
\usepackage{fancyhdr}
\usepackage{lastpage}
\pagestyle{fancy}
\fancyhf{}
\renewcommand{\headrulewidth}{0.6pt}
\lhead{\small 习近平新时代中国特色社会主义思想概论·Kimi·B卷}
\rhead{\small 共 \pageref{LastPage} 页}
\cfoot{\small 第 \thepage\ 页}

% ---- 计数器与样式 ----
\usepackage{enumitem}
\usepackage{xcolor}
\usepackage{tcolorbox}
\tcbuselibrary{skins, breakable}

% ---- 填空线 ----
\newcommand{\blank}[1]{\underline{\hspace{#1}}}

% ---- 答案盒子 ----
\newtcolorbox{answerbox}[1][]{
  breakable,
  colback=gray!6,
  colframe=gray!50,
  fonttitle=\bfseries\small,
  title=参考答案,
  #1
}

% ---- 题号格式 ----
\newcounter{bigq}
\newcommand{\BigQ}[1]{%
  \stepcounter{bigq}%
  \vspace{6pt}%
  \noindent{\large\bfseries 第\chinese{bigq}大题\quad #1}%
  \vspace{4pt}\par
}

\newcounter{subq}[bigq]
\newcommand{\SubQ}{%
  \stepcounter{subq}%
  \noindent\textbf{（\arabic{subq}）}\quad
}

% 分隔线
\newcommand{\examrule}{\noindent\rule{\linewidth}{0.6pt}}

% ---- 文档开始 ----
\begin{document}

% ============================================================
%  封面
% ============================================================
\begin{center}
  {\LARGE\bfseries 习近平新时代中国特色社会主义思想概论}\\[8pt]
  {\large\bfseries 期末全真模拟试卷（二）}\\[16pt]
  \begin{tabular}{lll}
    考试时间：120 分钟 & \quad & 满分：100 分 \\[4pt]
    \multicolumn{3}{l}{注意事项：请将答案写在答题纸指定区域，写在本卷上无效。}
  \end{tabular}
  \vspace{4pt}

  \examrule

  \vspace{6pt}
  \begin{tabular}{|c|c|c|c|c|c|}
    \hline
    \textbf{题号} & 一 & 二 & 三 & \textbf{总分} \\
    \hline
    \textbf{满分} & 30 & 20 & 50 & 100 \\
    \hline
    \textbf{得分} & & & & \\
    \hline
  \end{tabular}
\end{center}

\vspace{10pt}
\examrule
\vspace{6pt}

% ============================================================
%  第一大题：单项选择题
% ============================================================
\BigQ{单项选择题（每题 1 分，共 30 分）}

\SubQ "十四个坚持"在习近平新时代中国特色社会主义思想科学体系中的定位是：

A. 最为核心关键的组成部分，起支撑作用的"四梁八柱"

B. 新时代坚持和发展中国特色社会主义的基本方略，是落实这一思想的实践要求

C. 全面展示了这一思想的实践成果

D. 这一思想的世界观和方法论

\vspace{6pt}

\SubQ "六个必须坚持"在习近平新时代中国特色社会主义思想中的作用是：

A. 集中体现了这一思想的主要观点和基本主张

B. 构成了新思想的基本方略

C. 从世界观和方法论角度彰显了这一思想的理论品格与鲜明特征

D. 全面展示了这一思想的实践成果

\vspace{6pt}

\SubQ 关于"两个结合"的表述，以下哪项是正确的？

A. 马克思主义基本原理是"根脉"，中华优秀传统文化是"魂脉"

B. 马克思主义基本原理是"魂脉"，中华优秀传统文化是"根脉"

C. "两个结合"是指马克思主义基本原理同中国具体实际相结合、同世界优秀文化相结合

D. "两个结合"是习近平新时代中国特色社会主义思想的唯一理论来源

\vspace{6pt}

\SubQ 习近平新时代中国特色社会主义思想的历史地位不包括：

A. 当代中国马克思主义、二十一世纪马克思主义

B. 中华文化和中国精神的时代精华

C. 实现了马克思主义中国化时代化的"第三次飞跃"

D. 是党和国家必须长期坚持的指导思想

\vspace{6pt}

\SubQ "两个确立"与"两个维护"的关系是：

A. 两者完全相同，可以互相替代

B. "两个确立"是"两个维护"的政治前提，"两个维护"是"两个确立"的实践要求，二者紧密联系但内涵不同

C. "两个维护"包含"两个确立"，范围更广

D. 两者没有直接关联

\vspace{6pt}

\SubQ 中国特色社会主义是社会主义而不是其他什么主义，其根本依据在于：

A. 中国特色社会主义是在改革开放40多年的伟大实践中得来的

B. 中国特色社会主义既坚持了科学社会主义基本原则，又根据时代条件赋予其鲜明的中国特色

C. 中国特色社会主义是在对中华文明5000多年的传承发展中得来的

D. 中国特色社会主义是在近代以来中华民族由衰到盛180多年的历史进程中得来的

\vspace{6pt}

\SubQ 在"四个自信"中，以下哪项表述最准确？

A. 道路自信是对马克思主义理论科学性、真理性的自信

B. 理论自信是对发展方向和未来命运的自信

C. 制度自信是对中国特色社会主义制度具有制度优势的自信

D. 文化自信是对中国特色社会主义经济发展水平的自信

\vspace{6pt}

\SubQ 关于新时代伟大变革的里程碑意义，以下哪项不属于"三件大事"？

A. 迎来了中国共产党成立一百周年

B. 中国特色社会主义进入新时代

C. 完成脱贫攻坚、全面建成小康社会的历史任务，实现第一个百年奋斗目标

D. 全面实现社会主义现代化强国

\vspace{6pt}

\SubQ "五位一体"总体布局不包括以下哪项？

A. 经济建设

B. 政治建设

C. 国防建设

D. 生态文明建设

\vspace{6pt}

\SubQ "四个全面"战略布局中，以下哪项是新时代的新表述？

A. 全面建成小康社会

B. 全面建设社会主义现代化国家

C. 全面依法治国

D. 全面从严治党

\vspace{6pt}

\SubQ 全面建成社会主义现代化强国的战略安排中，从2020年到2035年的目标是：

A. 全面建成小康社会

B. 基本实现社会主义现代化

C. 建成富强民主文明和谐美丽的社会主义现代化强国

D. 实现全体人民共同富裕

\vspace{6pt}

\SubQ 中国式现代化的中国特色中，"全体人民共同富裕的现代化"强调的是：

A. 同步富裕、同时富裕

B. 允许一部分人先富起来，同时积极推动先富带后富

C. 平均主义分配

D. 只注重物质富裕

\vspace{6pt}

\SubQ 推进中国式现代化需要正确处理的重大关系中，不包括：

A. 顶层设计与实践探索的关系

B. 战略与策略的关系

C. 改革与革命的关系

D. 自立自强与对外开放的关系

\vspace{6pt}

\SubQ 团结奋斗是中国人民创造历史伟业的必由之路，这属于：

A. "四个全面"战略布局

B. "五个必由之路"之一

C. "十个明确"的内容

D. "十四个坚持"的内容

\vspace{6pt}

\SubQ 关于中国共产党的领导，以下表述错误的是：

A. 中国共产党领导是中国特色社会主义最本质的特征

B. 中国共产党领导是中国特色社会主义制度的最大优势

C. 中国共产党是最高政治领导力量

D. 党的领导是全面的、系统的、整体的，意味着党要事事具体包办

\vspace{6pt}

\SubQ 党的自身优势中，马克思主义政党第一位的能力是：

A. 思想引领力

B. 群众组织力

C. 政治领导力

D. 社会号召力

\vspace{6pt}

\SubQ 党中央集中统一领导是党的领导的：

A. 一般原则

B. 最高原则

C. 临时措施

D. 局部要求

\vspace{6pt}

\SubQ 关于人民立场，以下表述最准确的是：

A. 人民立场是中国共产党的基本政治立场

B. 人民立场是中国共产党的根本政治立场，是马克思主义政党区别于其他政党的显著标志

C. 人民立场只是口号，没有实际意义

D. 人民立场只适用于革命时期

\vspace{6pt}

\SubQ "民心是最大的政治，决定事业兴衰成败"这一论断说明：

A. 政治斗争是核心任务

B. 打江山、守江山，守的是人民的心

C. 民心可以忽视

D. 经济发展决定一切

\vspace{6pt}

\SubQ 关于共同富裕，以下表述错误的是：

A. 共同富裕是中国特色社会主义的本质要求

B. 共同富裕是中国式现代化的重要特征

C. 共同富裕是整齐划一、同步实现的

D. 促进共同富裕要把握好鼓励勤劳创新致富、坚持基本经济制度、尽力而为量力而行、坚持循序渐进的原则

\vspace{6pt}

\SubQ 群众路线是我们党的：

A. 一般工作方法

B. 生命线和根本工作路线

C. 临时性措施

D. 只适用于革命时期的路线

\vspace{6pt}

\SubQ 新时代全面深化改革开放就其性质而言是：

A. 一场深刻的革命

B. 一次简单的政策调整

C. 一项常规工作

D. 一个短期任务

\vspace{6pt}

\SubQ 全面深化改革必须牢牢把握总目标，实施好"六个紧紧围绕"的路线图，其中经济体制改革的重点是：

A. 紧紧围绕坚持党的领导、人民当家作主、依法治国有机统一

B. 紧紧围绕使市场在资源配置中起决定性作用和更好发挥政府作用

C. 紧紧围绕建设社会主义核心价值体系

D. 紧紧围绕提高科学执政、民主执政、依法执政水平

\vspace{6pt}

\SubQ 推进国家治理体系和治理能力现代化，必须：

A. 照搬西方政治制度

B. 坚定中国特色社会主义制度自信，把制度优势转化为国家治理效能

C. 放弃中国特色社会主义制度

D. 实行多党制

\vspace{6pt}

\SubQ 新发展理念中，"绿色发展"注重解决的是：

A. 发展动力问题

B. 发展不平衡问题

C. 人与自然和谐共生问题

D. 社会公平正义问题

\vspace{6pt}

\SubQ 构建新发展格局是：

A. 被迫之举和权宜之计

B. 封闭的国内单循环

C. 各地都搞自我小循环

D. 把握未来发展主动权的先手棋，是开放的国内国际双循环

\vspace{6pt}

\SubQ 社会主义基本经济制度中，"两个毫不动摇"是指：

A. 毫不动摇巩固和发展公有制经济，毫不动摇鼓励、支持、引导非公有制经济发展

B. 毫不动摇坚持计划经济，毫不动摇发展市场经济

C. 毫不动摇发展国有经济，毫不动摇限制私有经济

D. 毫不动摇坚持对外开放，毫不动摇保护国内市场

\vspace{6pt}

\SubQ 坚持教育优先发展、科技自立自强、人才引领驱动，加快建设教育强国、科技强国、人才强国，体现了：

A. 教育、科技、人才是全面建设社会主义现代化国家的一般性条件

B. 教育、科技、人才是全面建设社会主义现代化国家的基础性、战略性支撑

C. 只需要重视科技，教育和人才不重要

D. 只需要重视教育，科技和人才不重要

\vspace{6pt}

\SubQ 全过程人民民主是：

A. 过程民主和结果民主的统一

B. 过程民主和成果民主的统一

C. 只有形式没有实质的民主

D. 西方民主的翻版

\vspace{6pt}

\SubQ 关于意识形态工作，以下表述正确的是：

A. 意识形态工作无关紧要

B. 意识形态关乎旗帜、关乎道路、关乎国家政治安全，决定文化前进方向和发展道路

C. 意识形态可以放任自流

D. 意识形态只涉及文化领域

\vspace{10pt}
\examrule

% ============================================================
%  第二大题：多项选择题
% ============================================================
\BigQ{多项选择题（每题 2 分，共 20 分）}

\SubQ 习近平新时代中国特色社会主义思想是完整的科学体系，其内容包括：

A. "十个明确"

B. "十四个坚持"

C. "十三个方面成就"

D. "六个必须坚持"

E. "八个明确"

\vspace{6pt}

\SubQ 中国特色社会主义进入新时代，这个新时代的内涵包括：

A. 承前启后、继往开来、在新的历史条件下继续夺取中国特色社会主义伟大胜利的时代

B. 决胜全面建成小康社会、进而全面建设社会主义现代化强国的时代

C. 全国各族人民团结奋斗、不断创造美好生活、逐步实现全体人民共同富裕的时代

D. 全体中华儿女勠力同心、奋力实现中华民族伟大复兴中国梦的时代

E. 我国不断为人类作出更大贡献的时代

\vspace{6pt}

\SubQ 中国式现代化的中国特色包括：

A. 人口规模巨大的现代化

B. 全体人民共同富裕的现代化

C. 物质文明和精神文明相协调的现代化

D. 人与自然和谐共生的现代化

E. 走和平发展道路的现代化

\vspace{6pt}

\SubQ 推进中国式现代化必须牢牢把握的重大原则包括：

A. 坚持和加强党的全面领导

B. 坚持中国特色社会主义道路

C. 坚持以人民为中心的发展思想

D. 坚持深化改革开放

E. 坚持发扬斗争精神

\vspace{6pt}

\SubQ 健全党的全面领导制度，必须：

A. 完善党在各种组织中发挥领导作用的制度

B. 完善党协调各方的机制

C. 完善党领导各项事业的具体制度

D. 把党的领导贯彻到党和国家所有机构履行职责全过程

E. 发挥党在各种组织中总揽全局、协调各方的领导核心作用

\vspace{6pt}

\SubQ 坚持以人民为中心的发展思想，要求：

A. 始终把人民放在心中最高的位置

B. 想人民之所想，行人民之所嘱

C. 尊重人民主体地位，尊重人民首创精神

D. 人民是党的工作的最高裁决者和最终评判者

E. 让群众满意是党做好一切工作的价值取向和根本标准

\vspace{6pt}

\SubQ 全面深化改革开放要坚持的正确方法论包括：

A. 增强全面深化改革的系统性、整体性、协同性

B. 加强顶层设计和摸着石头过河相结合

C. 统筹改革发展稳定

D. 胆子要大，步子要稳

E. 坚持重大改革于法有据

\vspace{6pt}

\SubQ 新发展理念的实践要求包括：

A. 坚持创新在我国现代化建设全局中的核心地位

B. 不断增强发展的整体性

C. 推动形成绿色低碳的生产方式和生活方式

D. 推动形成更大范围、更宽领域、更深层次对外开放格局

E. 坚持全民共享、全面共享、共建共享、渐进共享

\vspace{6pt}

\SubQ 中国特色社会主义政治制度的显著优势包括：

A. 坚持党的集中统一领导，确保国家始终沿着社会主义方向前进

B. 坚持人民当家作主，紧紧依靠人民推动国家发展

C. 坚持全面依法治国，切实保障社会公平正义和人民权利

D. 坚持全国一盘棋，集中力量办大事

E. 坚持各民族一律平等，铸牢中华民族共同体意识

\vspace{6pt}

\SubQ 建设社会主义文化强国，必须：

A. 坚持马克思主义在意识形态领域指导地位的根本制度

B. 大力加强马克思主义理论建设

C. 积极塑造主流舆论新格局

D. 以社会主义核心价值观引领文化建设

E. 传承发展中华优秀传统文化，铸就社会主义文化新辉煌

\vspace{10pt}
\examrule

% ============================================================
%  第三大题：简答题
% ============================================================
\BigQ{简答题（每题 10 分，共 50 分）}

\SubQ（10 分）为什么说中国共产党的领导是中国特色社会主义制度的最大优势？请结合"中国最大的国情就是中国共产党的领导"和"健全党的全面领导制度"加以论述。

\vspace{6pt}

\SubQ（10 分）请阐述"江山就是人民，人民就是江山"的深刻内涵，并结合人民立场是中国共产党的根本政治立场，论述如何全面落实以人民为中心的发展思想。

\vspace{6pt}

\SubQ（10 分）为什么说新时代全面深化改革开放是一场深刻革命？请结合全面深化改革开放的正确方向加以论述。

\vspace{6pt}

\SubQ（10 分）请阐述人民民主是社会主义的生命，并结合全过程人民民主是社会主义民主政治的本质属性，论述如何坚定不移走中国特色社会主义政治发展道路。

\vspace{6pt}

\SubQ（10 分）为什么说文化繁荣兴盛是实现中华民族伟大复兴的必然要求？请结合建设具有强大凝聚力和引领力的社会主义意识形态、以社会主义核心价值观引领文化建设加以论述。

\vspace{10pt}
\examrule

\vspace{16pt}
\examrule
\vspace{4pt}
\begin{center}
  {\large\bfseries ——试题到此结束，请认真检查——}
\end{center}
\vspace{4pt}
\examrule

% ============================================================
%  参考答案
% ============================================================
\newpage
\setcounter{bigq}{0}

\begin{center}
  {\LARGE\bfseries 参考答案与评分标准}
\end{center}
\vspace{8pt}
\examrule

% -------- 第一大题参考答案 --------
\BigQ{单项选择题参考答案}

\begin{answerbox}

（1）B \quad （2）C \quad （3）B \quad （4）C \quad （5）B

（6）B \quad （7）C \quad （8）D \quad （9）C \quad （10）B

（11）B \quad （12）B \quad （13）C \quad （14）B \quad （15）D

（16）C \quad （17）B \quad （18）B \quad （19）B \quad （20）C

（21）B \quad （22）A \quad （23）B \quad （24）B \quad （25）C

（26）D \quad （27）A \quad （28）B \quad （29）B \quad （30）B

\end{answerbox}

% -------- 第二大题参考答案 --------
\BigQ{多项选择题参考答案}

\begin{answerbox}

（1）ABCD \quad （2）ABCDE \quad （3）ABCDE \quad （4）ABCDE \quad （5）ABCDE

（6）ABCDE \quad （7）ABCDE \quad （8）ABCDE \quad （9）ABCDE \quad （10）ABCDE

\end{answerbox}

% -------- 第三大题参考答案 --------
\BigQ{简答题参考答案}

\begin{answerbox}

\textbf{第1题（10分）：中国共产党的领导是中国特色社会主义制度的最大优势}

\textbf{一、中国最大的国情就是中国共产党的领导（3分）}

1. \textbf{中国共产党领导是中国特色社会主义最本质的特征}（1分）。这是由科学社会主义的理论逻辑、中国特色社会主义产生和发展的历史逻辑、中国特色社会主义迈向新征程的实践逻辑决定的。

2. \textbf{中国共产党领导是中国特色社会主义制度的最大优势}（1分）。党是中国特色社会主义制度的创建者，党的领导是中国特色社会主义制度优势发挥的根本保障，党的自身优势是中国特色社会主义制度优势的主要来源。

3. \textbf{中国共产党是最高政治领导力量}（1分）。这是马克思主义政党学说的基本原则，是对历史经验的深刻总结，是推进伟大事业的根本保证。

\textbf{二、党的自身优势是制度优势的主要来源（2分）}

党的自身优势包括：政治领导力（马克思主义政党第一位的能力）、思想引领力、群众组织力、社会号召力（1分）。这些优势通过党的领导制度体系转化为国家治理效能，使中国特色社会主义制度展现出强大生命力和巨大优越性（1分）。

\textbf{三、健全党的全面领导制度（3分）}

1. \textbf{党的领导制度是我国的根本领导制度}（1分）。它是党的领导核心地位的必然反映和内在要求，是中国特色社会主义制度体系的核心，是国家治理体系和治理能力现代化的关键。

2. 健全党的全面领导制度，必须\textbf{完善党在各种组织中发挥领导作用的制度}，完善党协调各方的机制，完善党领导各项事业的具体制度（1分）。

3. 必须把党的领导贯彻到党和国家\textbf{所有机构履行职责全过程}，发挥党总揽全局、协调各方的领导核心作用（1分）。

\textbf{四、加强党的全面领导为新时代党和国家事业发展提供了坚强保证（2分）}

党的十八大以来党和国家事业取得历史性成就、发生历史性变革，最根本的原因在于有习近平总书记作为党中央的核心、全党的核心掌舵领航，在于有习近平新时代中国特色社会主义思想科学指引（2分）。

\end{answerbox}

\vspace{6pt}

\begin{answerbox}

\textbf{第2题（10分）："江山就是人民，人民就是江山"与以人民为中心的发展思想}

\textbf{一、"江山就是人民，人民就是江山"的深刻内涵（4分）}

1. \textbf{人民是历史的创造者，是真正的英雄}（1分）。人民是创造历史的真正动力，是历史发展和社会进步的主体力量。

2. \textbf{打江山、守江山，守的是人民的心}（1分）。民心是最大的政治，决定事业兴衰成败。只有紧紧依靠最广大人民群众，为人民执好政、掌好权，才能始终赢得人民群众衷心拥护，实现长期执政。

3. \textbf{人民立场是中国共产党的根本政治立场}（1分）。这是马克思主义政党区别于其他政党的显著标志。坚持人民立场，就要始终牢记党的初心和使命，始终保持党同人民群众的血肉联系，热爱人民、尊重人民、敬畏人民。

4. \textbf{人民性是马克思主义的本质属性和鲜明品格}（1分）。坚持以人民为中心，是新时代坚持和发展中国特色社会主义的根本立场，是贯穿党治国理政全部活动的一条红线。

\textbf{二、全面落实以人民为中心的发展思想（6分）}

1. \textbf{坚持和贯彻党的群众路线}（2分）。群众路线是我们党的生命线和根本工作路线，是我们党永葆青春活力和战斗力的重要传家宝。贯彻党的群众路线，"知"是基础、是前提，"行"是重点、是关键。调查研究是获得真知灼见的源头活水，是贯彻群众路线的有效途径。

2. \textbf{把为人民造福的事情真正办好办实}（2分）。要落实到新时代中国特色社会主义的各项事业、全部工作之中，着力解决好人民群众最关心最直接最现实的利益问题。要正确把握民生和发展的关系，坚守底线、突出重点、完善制度、引导预期，坚持尽力而为、量力而行。

3. \textbf{推动全体人民共同富裕取得更为明显的实质性进展}（2分）。共同富裕是中国特色社会主义的本质要求，是中国式现代化的重要特征。要从全局角度来把握共同富裕，处理好先富和共富的关系，坚持正确的原则和科学的思路，在高质量发展中促进共同富裕，形成中间大、两头小的橄榄型分配结构。

\end{answerbox}

\vspace{6pt}

\begin{answerbox}

\textbf{第3题（10分）：新时代全面深化改革开放是一场深刻革命}

\textbf{一、新时代全面深化改革开放是一场深刻革命的原因（5分）}

1. \textbf{改革开放是我们前进的重要法宝}（1分）。改革开放是党和人民大踏步赶上时代的重要法宝，是坚持和发展中国特色社会主义的必由之路。新时代坚持和发展中国特色社会主义，实现"两个一百年"奋斗目标、实现中华民族伟大复兴，必须以改革开放为强大动力。

2. \textbf{新时代全面深化改革开放就其艰巨性、复杂性和系统性来说，是一场深刻的革命}（2分）。这是涉险滩、闯难关的变革，是思想理论的深刻变革、改革组织方式的深刻变革、国家制度和治理体系的深刻变革、人民广泛参与的深刻变革。

3. \textbf{新时代的改革开放是全方位、深层次、根本性的}（2分）。是在多年改革开放基础上的深化，是一场全面、系统、整体的制度创新，使社会主义制度的优越性得到更好发挥，取得的成就是历史性、革命性、开创性的。

\textbf{二、坚持全面深化改革开放的正确方向（5分）}

1. \textbf{必须坚持和改善党的全面领导、坚持和完善中国特色社会主义制度}（2分）。全面深化改革开放既不走封闭僵化的老路，也不走改旗易帜的邪路，必须以是否有利于坚持和改善党的全面领导、坚持和完善中国特色社会主义制度为根本标准。

2. \textbf{必须坚持以人民为中心，促进社会公平正义、增进人民福祉}（1分）。社会主义改革开放的出发点和落脚点，是为了更好实现和维护人民利益、为了让老百姓过上好日子。

3. \textbf{必须有利于进一步解放思想、进一步解放和发展社会生产力、进一步解放和增强社会活力}（2分）。习近平指出："这'三个进一步解放'既是改革的目的，又是改革的条件。"解放思想是思想前提，解放和发展社会生产力是鲜明特征和首要任务，解放和增强社会活力是内在要求和基本目的。

\end{answerbox}

\vspace{6pt}

\begin{answerbox}

\textbf{第4题（10分）：人民民主是社会主义的生命与中国特色社会主义政治发展道路}

\textbf{一、人民民主是社会主义的生命（3分）}

1. \textbf{民主是全人类共同价值}，是人类政治文明发展的成果（1分）。民主是历史的、具体的，一个国家的民主制度是在这个国家的社会历史条件下形成的，植根于本国的历史文化传统。

2. \textbf{人民民主建立在社会主义经济基础之上}，体现了社会主义国家的性质，反映了社会主义制度的本质要求，是一种新型的社会主义民主（1分）。

3. \textbf{人民民主是全面建设社会主义现代化国家的应有之义}（1分）。社会主义现代化是全体人民的现代化，必然要求具有高度发达的民主政治，保障人民当家作主，不断推动人的全面发展。

\textbf{二、全过程人民民主是社会主义民主政治的本质属性（4分）}

1. \textbf{全过程人民民主是社会主义民主政治的伟大创造}（1分）。它植根于我国社会主义民主的长期实践，是中国特色社会主义民主政治理论和实践的重大成果。

2. \textbf{全过程人民民主是全链条、全方位、全覆盖的民主}（1分）。包括民主选举、民主协商、民主决策、民主管理、民主监督各个环节，涵盖经济、政治、文化、社会、生态文明等各个方面。

3. \textbf{全过程人民民主体现"四个相统一"}（1分）：过程民主和成果民主的统一、程序民主和实质民主的统一、直接民主和间接民主的统一、人民民主和国家意志的统一。

4. \textbf{全过程人民民主是最广泛、最真实、最管用的民主}（1分）。既能确保人民充分享有各方面权利，又能有效促进国家治理高效和社会和谐稳定，维护和保障人民群众的根本利益。

\textbf{三、坚定不移走中国特色社会主义政治发展道路（3分）}

1. \textbf{必须坚持党的领导、人民当家作主、依法治国有机统一}（1分）。党的领导是人民当家作主和依法治国的根本保证，人民当家作主是社会主义民主政治的本质特征，依法治国是党领导人民治理国家的基本方式。

2. \textbf{必须积极稳妥推进政治体制改革}（1分）。

3. \textbf{必须始终保持政治定力}（1分），坚定对中国特色社会主义政治制度的自信，增强走中国特色社会主义政治发展道路的信心和决心。

\end{answerbox}

\vspace{6pt}

\begin{answerbox}

\textbf{第5题（10分）：文化繁荣兴盛与社会主义文化强国建设}

\textbf{一、文化繁荣兴盛是实现中华民族伟大复兴的必然要求（4分）}

1. \textbf{文化繁荣兴盛是实现中华民族伟大复兴的精神支撑}（1分）。没有高度的文化自信，没有文化的繁荣兴盛，就没有中华民族伟大复兴。

2. \textbf{文化繁荣兴盛是建设社会主义现代化强国的应有之义}（1分）。社会主义现代化强国是物质文明和精神文明相协调的现代化。

3. \textbf{文化繁荣兴盛是满足人民日益增长的美好生活需要的内在要求}（1分）。人民对精神文化生活的需求日益增长，需要更加丰富、更高质量的文化供给。

4. \textbf{文化繁荣兴盛是在世界文化激荡中站稳脚跟的前提基础}（1分）。只有坚定文化自信，才能在世界文化激荡中保持自身文化独立性和主体性。

\textbf{二、建设具有强大凝聚力和引领力的社会主义意识形态（3分）}

1. \textbf{意识形态关乎旗帜、关乎道路、关乎国家政治安全}，决定文化前进方向和发展道路（1分）。

2. \textbf{坚持马克思主义在意识形态领域指导地位的根本制度}（1分）。这是坚持和加强党对宣传文化事业全面领导的本质要求，是恪守党的本质属性、巩固党的团结统一的必然要求，是坚持正确发展道路、实现国家长治久安的必然要求。

3. \textbf{大力加强马克思主义理论建设，积极塑造主流舆论新格局}（1分）。把马克思主义指导地位贯穿到文化建设各方面，实施马克思主义理论研究和建设工程，落实意识形态工作责任制。

\textbf{三、以社会主义核心价值观引领文化建设（3分）}

1. \textbf{广泛践行社会主义核心价值观}（1分）。社会主义核心价值观是当代中国精神的集中体现，凝结着全体人民共同的价值追求。要使其成为广泛的社会共识，内化为人们的精神追求，外化为人们的自觉行动。

2. \textbf{弘扬以伟大建党精神为源头的中国共产党人精神谱系}（1分）。伟大建党精神是中国共产党的精神之源，是激励全党全国各族人民奋勇前进的强大精神动力。

3. \textbf{提高全社会文明程度，铸就社会主义文化新辉煌}（1分）。传承发展中华优秀传统文化，推动其创造性转化、创新性发展；繁荣发展文化事业和文化产业；不断提升国家文化软实力和中华文化影响力，加强国际传播能力建设。

\end{answerbox}

\vspace{12pt}
\examrule
\begin{center}
  \textit{参考答案仅供参考，评分时应酌情给分，观点正确、论述充分即可得分。}
\end{center}

\end{document}